\documentclass[xcolor=dvipsnames]{beamer} % dvipsnames gives more built-in colors
\usetheme{Singapore}
\useoutertheme{miniframes} % Alternatively: miniframes, infolines, split
\useinnertheme{circles}
\usecolortheme{dove}

\newcommand{\BE}{ \begin{enumerate}}
\newcommand{\EE}{ \end{enumerate}}
\newcommand{\BI}{ \begin{itemize}}
\newcommand{\EI}{ \end{itemize}}
\newcommand{\BEQ}{ \begin{equation}}
\newcommand{\EEQ}{ \end{equation}}
\newcommand{\BDM}{ \begin{displaymath}}
\newcommand{\EDM}{ \end{displaymath}}
\newcommand{\BFNS}{ \begin{footnotesize}}
\newcommand{\EFNS}{ \end{footnotesize}}
\newcommand{\BS}{ \begin{small}}
\newcommand{\ES}{ \end{small}}
\newcommand{\BSS}{ \begin{scriptsize}}
\newcommand{\ESS}{ \end{scriptsize}}
\newcommand{\BTi}{ \begin{tiny}}
\newcommand{\ETi}{ \end{tiny}}
\newcommand{\BVTi}{ \begin{verytiny}}
\newcommand{\EVTi}{ \end{verytiny}}
\newcommand{\balpha}{\mbox{\boldmath $\alpha$}}
\newcommand{\bOmega}{\mbox{\boldmath $\Omega$}}
\newcommand{\btau}{\mbox{\boldmath $\tau$}}
\newcommand{\bchi}{\mbox{\boldmath $\chi$}}
\newcommand{\bnabla}{\mbox{\boldmath $\nabla$}}
\newcommand{\grad}{\mbox{\boldmath $\nabla$}}
\newcommand{\boldeta}{\mbox{\boldmath $\eta$}}
\newcommand{\bphi}{\mbox{\boldmath $\varphi$}}
\newcommand{\ds}{\displaystyle}
\newcommand{\nl}{\ \\ }
\newcommand{\ud}{\textrm{ d}}
\newcommand{\bs}{\bigskip}
\newcommand{\BTc}{ \begin{tabular}{c}}
\newcommand{\BTl}{ \begin{tabular}{l}}
\newcommand{\BTr}{ \begin{tabular}{r}}
\newcommand{\ET}{ \end{tabular}}
\newcommand{\bu}{\mathbf{u}}
\newcommand{\bv}{\mathbf{v}}
\newcommand{\bx}{\mathbf{x}}
\newcommand{\be}{\mathbf{e}}
\newcommand{\bb}{\mathbf{b}}
\newcommand{\bk}{\mathbf{k}}
\newcommand{\bd}{\mathbf{d}}
\newcommand{\bn}{\mathbf{n}}
\newcommand{\bq}{\mathbf{q}}
\newcommand{\ba}{\mathbf{a}}
\newcommand{\bF}{\mathbf{F}}
\newcommand{\bI}{\mathbf{I}}
\newcommand{\bT}{\mathbf{T}}
\newcommand{\ust}{u_{*}}
\newcommand{\ddz}{\frac{\partial}{\partial z}}
\newcommand{\ddt}{\frac{\partial}{\partial t}}
\newcommand{\dudx}{\frac{\partial u}{\partial x}}
\newcommand{\dudz}{\frac{\partial u}{\partial z}}
\newcommand{\dvdz}{\frac{\partial v}{\partial z}}
\newcommand{\dudt}{\frac{\partial u}{\partial t}}
\newcommand{\dvdt}{\frac{\partial v}{\partial t}}
\newcommand{\ddudxx}{\frac{\partial^2 u}{\partial x^2}}
\newcommand{\dbudz}{\frac{\partial \bu}{\partial z}}
\newcommand{\dbudt}{\frac{\partial \bu}{\partial t}}
\newcommand{\partd}[2]{\frac{\partial #1}{\partial #2}}
\newcommand{\lagd}[2]{\frac{D #1}{D #2}}
\newcommand{\lagda}[2]{\frac{D_a #1}{D #2}}

\newcommand{\arrow}[5]{ \put(#1,#2){\vector(#3,#4){#5}} }
\newcommand{\grid}{ 
\put(0,0){\vector(1,0){8}}
\put(0,1){\line(1,0){7}}  \put(0,2){\line(1,0){7}}
\put(0,3){\line(1,0){7}}  \put(0,4){\line(1,0){7}}
\put(0,5){\line(1,0){7}}
\put(0,0){\vector(0,1){6}}  
\put(1,0){\line(0,1){5}}  \put(2,0){\line(0,1){5}}  
\put(3,0){\line(0,1){5}}  \put(4,0){\line(0,1){5}}  
\put(5,0){\line(0,1){5}}  \put(6,0){\line(0,1){5}}  
\put(7,0){\line(0,1){5}}
\put(2.8,-0.7){$j$}
\put(7.3,-0.7){$x$}
\put(-0.7,2.8){$n$}
\put(-2.0,3.8){$n+1$}
\put(-0.7,5.3){$t$}}
\newcommand{\arakawa}{
\put(1,0){\line(0,1){4}}
\put(3,0){\line(0,1){4}}
\put(0,1){\line(1,0){4}}
\put(0,3){\line(1,0){4}}
}


\title{Numerical modelling of the atmosphere and oceans \\ Lecture 4}
\author{P.L. Vidale \footnote{\tiny{based in part on J. Methven lecture notes from 2013}}}
\institute{
  National Centre for Atmospheric Science and Department of Meteorology\\
  University of Reading, UK\\
  p.l.vidale@reading.ac.uk}
\date{Spring Term 2023}

\begin{document}

\begin{frame}
	\titlepage
\end{frame}

\section[Outline]{}
\begin{frame}
	\tableofcontents
\end{frame}

\section{Fluids on rotating Earth}
\subsection{}

\frame { 
	\frametitle{From Newton to Navier-Stokes}

\BSS

 In fluid dynamics we often thing of Newton's 2nd law in this form:
 
 
 \begin{equation}
 	\lagda{\mathbf{v_a}}{t} = \sum \mathbf{F}
 \end{equation}
 where the subscript a indicates that we follow the motion as viewed in an inertial system. In fact, Newton's 2nd law and the continuity equation (of mass
 and entropy, with appropriate {\em boundary conditions} and {\em forcing}) are often presented together, as the Navier-Stokes equations:

\begin{figure}[H]
	\includegraphics[trim={0cm 0.8cm 0cm 1cm},clip,width=6.cm]{/Users/vidale/Projects/Teaching/MTMW14/Figures/Navier-Stokes}
	\label{fig:Navier-Stokes}
	\setbeamerfont{caption name}{size=\tiny}
	\caption{\label{fig:blue_rectangle} \tiny The Navier-Stokes equations}
\end{figure}

which are capable of representing any type of fluid motion, and are appicable to an extremely wide range of scales. 

We do not know how to solve these equations, and in fact the Clay Mathematics Institute of Cambridge, Massachusetts (CMI) selected this as one of the seven Millenium Prize Problems, offering \$ 1 million for a solution.

\ESS }



\frame { 
	\BSS
We need to relate Newton's 2nd law, as viewed in an inertial system,

\begin{equation}
	\lagda{\mathbf{v_a}}{t} = \sum \mathbf{F}
\end{equation}


 to a rotating coordinate system point of view. This requires expressing $v_a$ as $v$ by using the relationship for a position vector $r$.

\begin{equation}
	\lagda{\mathbf{r}}{t} = \lagd{\mathbf{r}}{t} + \Omega \times \mathbf{r}
\end{equation}

and $\mathbf{v_a} = \mathbf{v} + \Omega \times \mathbf{r}$, so that:

\begin{equation}
	\lagda{\mathbf{v_a}}{t} = \lagd{\mathbf{v_a}}{t} + \Omega \times \mathbf{v_a}
\end{equation}

and finally:

\begin{equation}
	\lagda{\mathbf{v_a}}{t} = \lagd{\mathbf{v}}{t} + 2 \Omega \times \mathbf{v} - \Omega^2 \mathbf{r}
\end{equation}

\ESS }

\begin{frame}  
\frametitle{Still Newton's 2nd law, now for fluid motion on Earth}
\BSS
The form of the momentum equations for geophysical fluids is:
\begin{equation}
	\lagd{\mathbf{v}}{t} =-\frac {\nabla p}{\rho}+\mathbf{g}+\mathbf{F_r}
\end{equation}

or, if Earth's rotation is included:

\begin{equation}
	\lagd{\mathbf{v}}{t} =-\frac {\nabla p}{\rho}+\mathbf{g}+\mathbf{F_r} - 2 \Omega \times \mathbf{v}
\end{equation}
~

where $F_r$ now represents frictional forces. The system supports a complex variety of phenomena across a huge range of scales.

Flow almost adiabatic with ``weak'' forcing $\Rightarrow$ stable
stratification established.

Earth's rotation has profound influence on motions $\Rightarrow$
layerwise-2D motions.

~

The system may be simplified by noting the following:

\begin{itemize}
\item
Earth rotates about its polar axis at rate $\Omega=7.292 \times 10^{-5}$s$^{-1}$,

\item
Earth is approximately an oblate spheroid,

\item 
Radius to poles $a_{pole}\approx 6357$km, radius to equator $a_{eq}\approx
6378$km. Average, $a\approx 6371$km.

\item
Atmosphere and oceans are shallow relative to the Earth's radius.
\end{itemize}
\ESS

\end{frame}

\frame { \frametitle{Shallow atmosphere approximation} 

\BSS Also known as the {\em traditional approximation} since it also
applies to the oceans. The deepest motions in the atmosphere have
depth scales, $H\approx\,$20km. The oceans are even shallower. Therefore
$H/a< 1/300$. The shallow limit $H/a\ll 1$ gives:

\begin{itemize}
\item
radial coordinate $r\rightarrow a$ (i.e., constant)

\item
$\partd{}{r}\rightarrow \partd{}{z}$ where $z$ is height relative to Earth's surface

\item
In order to retain same form for kinetic energy and zonal angular
momentum equations, must drop terms containing $w$ in horizontal
momentum equations and the horizontal component of Coriolis
acceleration.
\end{itemize}

The resulting {\em primitive equations} (PEs) are the basis of almost
all atmosphere and ocean {\em general circulation models} (GCMs)
except the Met Office Unified Model which does not make a shallow
atmosphere approximation.

\begin{itemize}
\item
Biggest errors are in Tropics associated with neglect of horizontal
component of Coriolis acceleration.\footnote{\BTi White, Hoskins, Roulstone and Staniforth (2005) Consistent approximate models of the global atmosphere: shallow, deep, hydrostatic, quasi-hydrostatic and non-hydrostatic. \emph{Quart. J. Roy. Met. Soc.}, \textbf{131},
2081-2107.
\ETi }

\item
Effects of spherical geopotential approximation are unknown.\footnote{\BTi White, Staniforth and Wood (2008) Spheroidal coordinate systems for modelling global atmospheres. \emph{Quart. J. Roy. Met. Soc.}, \textbf{134},
261-270.
\ETi }
\end{itemize}

\ESS }

\frame { \frametitle{Filtering the equations to remove unwanted scales of motion} 
	
\BSS

The problem with our Euler equations is that they support all types of motion, some of which may be unwanted, depending on what scientific questions we are trying to answer. 

\medskip

{\bf Example: sound waves.}\\
The governing equations of compressible fluid dynamics contain acoustic waves. In low Mach number flows, typical of atmospheric or oceanic conditions, acoustic waves play no essential role and yet they severely constrain the time step in numerical modeling. Thus, there has long been an interest in approximating the governing equations to eliminate or "filter out" acoustic waves. 

\medskip

\begin{tabular}{|c|c|c|c|c|}
	\hline
	Wave & horizontal & propagation & propagation  & time step \\
	type & dimension & speed & direction & required \\
	\hline
	Sound &  &  &  & \\
	\hline
	Short Gravity &  &  &  & \\
	\hline
	Long Gravity &  &  &  & \\
    \hline
	Rossby &  &  & & \\
	\hline
	Kelvin &  &  &  & \\
	\hline
\end{tabular}

\ESS }

\frame { \frametitle{Further approximations} 

\BSS Further approximations are usually based on {\em scaling} where
the typical spatial and length scales associated with motions of
interest are assumed. For example:

~

{\bf Planar approximation ($L/a \ll 1$)} Can use local Cartesian coordinates where $(dx,dy)=(a \cos\phi_0
\,d\lambda, a\,d\phi)$. 

Can drop spherical metric terms from PEs.

~

{\bf Anelastic approximation ($\Delta \rho /\rho_0 \ll 1$)} Mass conservation equation becomes $\nabla.(\rho_r {\mathbf u})=0$
where $\rho_r(z)$ is a reference density.
Filters sound waves.

~

{\bf Boussinesq approximation ($\Delta \rho /\rho_0 \ll 1$ and $H\ll H_{\rho}$)} Huge density height scale, $H_{\rho}$, implies $w \partd{\rho_r}{z}\ll
\rho_r \partd{w}{z}$. In practice, the density variation is only important in the buoyancy term of the (vertical) momentum equation, $\rho g$. The continuity equation:
\begin{equation}
	\frac{1}{\rho} \lagd{\rho}{t} + \nabla \cdot \mathbf{u} = 0
\end{equation}

ends up simplified to its incompressible ($\rho_r \approx$ constant) form: $\nabla \cdot \mathbf{u} = 0$. 

~

{\bf Hydrostatic balance ($H/L \ll 1$)} Vertical momentum equation reduces to:
\begin{equation}
\frac{1}{\rho}\partd{p}{z}=-g
\label{hydrostatic}
\end{equation}
Filters vertically-propagating sound waves, but can support Lamb waves, and can distort gravity waves if $L_x \approx L_y$.

If we also want to filter out gravity waves, we must additionally set the local time rate of change of the divergence field to zero.

\ESS }


\section{Simplifying the PEs}
\subsection{}
\frame { \frametitle{PEs on the sphere}
\BSS
In an Eulerian framework, our equation (1) ends up looking like this:
\begin{eqnarray}
\frac {\partial u}{\partial t} + \left(\mathbf{v}\cdot\nabla\right)u=-\frac{1}{\rho} \frac{\partial p}{\partial x}+F_x \\
\frac {\partial v}{\partial t} + \left(\mathbf{v}\cdot\nabla\right)v=-\frac{1}{\rho} \frac{\partial p}{\partial y}+F_y \\
\frac {\partial w}{\partial t} + \left(\mathbf{v}\cdot\nabla\right)w=-\frac{1}{\rho} \frac{\partial p}{\partial z}-g+F_z
\end{eqnarray}

We must now consider the fact that we are studying fluids (atmosphere, ocean) on our planet, which is approximately a sphere of radius $a$, rotating with angular velocity $\Omega$,

\begin{minipage}{0.4\textwidth}
	\begin{figure}[H]
		\includegraphics[trim={0cm 2cm 0cm 2cm},clip,width=3.cm]{/Users/vidale/Projects/Teaching/MTMW14/"2016 module"/RotatingEarthDiagram.pdf}
		\label{fig:Spherical}
		\setbeamerfont{caption name}{size=\tiny}
		\caption{\label{fig:blue_rectangle} \tiny Schematic of spherical coordinate system}
	\end{figure}
\end{minipage} \hfill
\begin{minipage}{0.55\textwidth}
In the spherical coordinate system we use $\lambda$ (longitude), $\theta$ (latitude) and r, to replace $x,y,z$, where $x=r \cos{\theta} \lambda$; $y=r\theta$; $z=r-a$, so that:
\begin{eqnarray}
 u=\frac {dx}{dt}=r \cos{\theta} \frac {d \lambda}{d t} ; v=\frac {dy}{dt}= r\frac {d\theta}{dt}= ; w=\frac {dz }{dt }
\end{eqnarray}
are how we write our spatial derivatives on the Sphere.	
\end{minipage}
\ESS }

\frame { \frametitle{PEs on the Sphere before scale analysis}
	\BSS
	\includegraphics[trim={2cm 10.5cm 0cm 2.cm},clip,width=10.cm]{/Users/vidale/Projects/Teaching/MTMW14/"2016 module"/PEs_Sphere_1.pdf}
	\label{fig:Spherical-PEs1}
	\setbeamerfont{caption name}{size=\tiny}
	%\caption{\label{fig:blue_rectangle} \tiny PEs on the Sphere}
	
		\includegraphics[trim={2cm 5cm 0cm 4cm},clip,width=10.cm]{/Users/vidale/Projects/Teaching/MTMW14/"2016 module"/PEs_Sphere_2.pdf}
	\label{fig:Spherical-PEs2}
	\setbeamerfont{caption name}{size=\tiny}
	%\caption{\label{fig:blue_rectangle} \tiny PEs on the Sphere}
	
	\ESS }


\frame { \frametitle{Small quantities, big quantities: Scale Analysis}
\BSS

Cancelling out the smallest terms, we end up with simpler equations and some important approximations, for instance hydrostatic in the vertical momentum equation.

\begin{figure}[H]
	\includegraphics[trim={0cm 0cm 0cm 1cm},clip,width=11.cm]{/Users/vidale/Projects/Teaching/MTMW14/"2016 module"/ScaleAnalysisTables.pdf}
	\label{fig:Spherical}
	\setbeamerfont{caption name}{size=\tiny}
	\caption{\label{fig:blue_rectangle} \tiny Scale Analysis Tables}
\end{figure}
\ESS }

\frame { \frametitle{Simplified PEs after scale analysis}
\BSS
		\includegraphics[trim={2cm 2.5cm 0cm 2.cm},clip,width=12.cm]{/Users/vidale/Projects/Teaching/MTMW14/"2016 module"/PEs_Sphere_3.pdf}
	\label{fig:Spherical-PEs1}
	\setbeamerfont{caption name}{size=\tiny}
	%\caption{\label{fig:blue_rectangle} \tiny PEs on the Sphere}
	
\ESS }

\frame { \frametitle{Simplified PEs after scale analysis, now back on a PLANE}
\BSS
$f$-plane: Coriolis term is a constant $f_0$ over the entire domain\\
$\beta$-plane: Coriolis term changes linearly in the meridional direction\\

	\begin{eqnarray}
		\partd{u}{t} + \left( \mathbf{v} \cdot \nabla\right) u - f_0 v & = & -\frac{1}{\rho_o}\partd{p}{x} +F_x \\
		\partd{v}{t} + \left( \mathbf{v} \cdot \nabla\right) v + f_0 u & = & -\frac{1}{\rho_o}\partd{p}{y} +F_y \\
		 0 & = & - \frac{1}{\rho} \partd{p}{z} - g + F_z
		%\partd{u}{x}+\partd{v}{y}+\partd{w}{z} & = & 0
	\end{eqnarray}

and the continuity equation also becomes much simpler, as metric terms drop out. The simpler geometry of the plane may be retained while taking the latitudinal varations of $f$ into account, by defining a $\beta$-plane, for which $f$ is calculated by linearising the Coriolis terms around a constant $f_0=2\Omega \sin\phi_0$ -
the component of the Earth's rotation normal to the Earth's surface at
reference latitude $\phi_0$. 	
		
%\includegraphics[trim={2cm 0.5cm 0.cm 10.cm},clip,width=12.cm]{/Users/vidale/Projects/Teaching/MTMW14/"2016 module"/PEs_Sphere_4.pdf}
%\label{fig:Spherical-PEs1}
%\setbeamerfont{caption name}{size=\tiny}
%\caption{\label{fig:blue_rectangle} \tiny PEs on the Sphere}

\ESS }

\section{Shallow water equations}
\subsection{}
\frame { \frametitle{Shallow water equations (I)} 

\BSS Making the planar, Boussinesq and hydrostatic approximations, and neglecting those unresolved forces, 
the PEs reduce to:
\begin{eqnarray}
\lagd{u}{t}-f_0 v & = & -\frac{1}{\rho_o}\partd{p}{x} \\
\lagd{v}{t}+f_0 u & = & -\frac{1}{\rho_o}\partd{p}{y} \\
		 0 & = & - \frac{1}{\rho} \partd{p}{z} - g \\
\partd{u}{x}+\partd{v}{y}+\partd{w}{z} & = & 0
\end{eqnarray}

plus thermodynamic eqn and eqn of state. Remember that $f_0=2\Omega \sin\phi_0$ is the Coriolis parameter normal to the Earth's surface at
reference latitude $\phi_0$.

~

Pressure and fluid depth are next divided into
a reference state and perturbation:
\begin{equation*}
p=p_r(z)+p'(x,y,z,t)~~~~~;~~~~~h=H+\eta(x,y,z,t)
\end{equation*}

\ESS }

\frame { \frametitle{Shallow water equations (II)} 

\BSS The shallow water equations are obtained by integrating the PEs
over a fluid layer of depth $h$. Integrating hydrostatic balance
downwards from the top to any level $z$:
\begin{eqnarray*}
\frac{1}{\rho}\partd{p}{z}=-g \Rightarrow p_{top}-p(x,y,z,t) & = & -\rho_o g (h-z) \\
\Rightarrow -\frac{1}{\rho_o} \partd{p}{x} & = & -g \partd{h}{x}
\end{eqnarray*}

assuming that there are no pressure perturbations on the top.
The pressure gradient terms become $-g\partd{\eta}{x}$ and
$-g\partd{\eta}{y}$, since the reference depth $H$ must be constant.

The mass conservation equation integrated vertically is:
\begin{eqnarray*}
\partd{u}{x}+\partd{v}{y}+\partd{w}{z} = 0 \Rightarrow \left\{ \partd{u}{x}+\partd{v}{y} \right\} h+[w]^{top}_{bot} & = & 0 \\
\left\{ \partd{u}{x}+\partd{v}{y} \right\} h+\lagd{h}{t} & = & 0 \\
\partd{h}{t}+\partd{(hu)}{x}+\partd{(hv)}{y} & = & 0
\end{eqnarray*}
where $(u,v)$ is now the depth-average velocity and $w_{bot}=0$ was assumed. 

\ESS }


\frame { \frametitle{SWEs linearised around a basic state} 
	
	\BSS {\bf Shallow-water equations} on the
	plane tangent to the Earth's surface, {\em linearised} about a resting basic state ($u=u_r(z)+u'(x,y,z,t); v=v_r(z)+v'(x,y,z,t);h=H+\eta(x,y,z,t)$):
	\begin{eqnarray}
	&&\frac{\partial \eta}{\partial t} + H \left(
	\frac{\partial u}{\partial x}+\frac{\partial v}{\partial y}\right) = 0,
	\\
	&&\frac{\partial u}{\partial t} - f_0 v = - g\frac{\partial \eta}{\partial x}, \\
	&&\frac{\partial v}{\partial t} + f_0 u = - g\frac{\partial \eta}{\partial y}, 
	\end{eqnarray}
	where $\eta$ is the surface elevation, $H$ is fluid depth, $(u,v)$ is
	the depth-averaged horizontal velocity, $f$ is the Coriolis parameter
	and $g$ is the gravitational acceleration. 
	
	~~
	
	Now discretise the problem using regular grids to describe $u$, $v$ and $\eta$ 
	
	e.g., $\eta_{ij}=$ free surface elevation at the $i$th longitude and $j$th
	latitude point.

\ESS }

\frame { \frametitle{Horizontal gradients in SWEs: how to discretize in space?} 
\BSS 	
	SWEs seem nice and simple, but now we have spatial gradients: $\frac{\partial u}{\partial x}$, $\frac{\partial v}{\partial y}$, $\frac{\partial \eta}{\partial x}$, $\frac{\partial \eta}{\partial y}$.\\

~
	
	How are we going to discretise the problem, using regular grids, to obtain a good simulation of $u$, $v$ and $\eta$? Something like this?
	
	\begin{tabular}{lc}
		\begin{minipage}[c]{0.7\textwidth}
			\begin{eqnarray*}
				\frac{\partial u}{\partial x}\Big|_{ij} &\approx& \left(\frac{u_{i+1,j}+u_{i+1,j-1}}{2} - \frac{u_{i-1,j}+u_{i-1,j-1}}{2} \right)\frac{1}{2 \Delta x} \\
				\frac{\partial v}{\partial y}\Big|_{ij} &\approx& \left(\frac{v_{i,j+1}+v_{i-1,j+1}}{2} - \frac{v_{i,j-1}+v_{i-1,j-1}}{2} \right)\frac{1}{2 \Delta y} \\
				&& \\
				\frac{\partial \eta}{\partial x}\Big|_{ij} &\approx& \left(\frac{\eta_{i+1,j}+\eta_{i+1,j+1}}{2} - \frac{\eta_{i-1,j}+\eta_{i-1,j+1}}{2} \right)\frac{1}{2 \Delta x} \\
				f v\big|_{ij} &\approx& f_j \frac{v_{i,j}+v_{i,j+1}+v_{i-1,j}+v_{i-1,j+1}}{4} \\
				&& \\
				\frac{\partial \eta}{\partial y}\Big|_{ij} &\approx& \left(\frac{\eta_{i,j+1}+\eta_{i+1,j+1}}{2} - \frac{\eta_{i,j-1}+\eta_{i+1,j-1}}{2} \right)\frac{1}{2 \Delta y} \\
				f u\big|_{ij} &\approx& f_j \frac{u_{i,j}+u_{i+1,j}+u_{i,j-1}+u_{i+1,j-1}}{4}
			\end{eqnarray*}
		\end{minipage}
\end{tabular}

~

	{\bf But why? It all seems so arbitrary.} There are a few criteria: we want to avoid computational modes and we want waves to propagate/evolve in a realistic way in space and time.

\ESS
}

\frame  {
	\frametitle{A preview of Lecture 5: graphical representation of inertia-gravity wave dispersion ($R_D/d = 10$), fine grid}
 	
	\begin{center}
		\begin{tabular}{cc}
			\includegraphics[width=0.45\textwidth]{Figures/gravity_exact_dispersion_Rd_10.eps} &
			\includegraphics[width=0.45\textwidth]{Figures/gravity_Agrid_dispersion_Rd_10.eps}
			\\
			\includegraphics[width=0.45\textwidth]{Figures/gravity_Bgrid_dispersion_Rd_10.eps} &
			\includegraphics[width=0.45\textwidth]{Figures/gravity_Cgrid_dispersion_Rd_10.eps}
		\end{tabular}
	\end{center}
	\BSS Shading shows $\omega/f_0$ ranging from blue (low) to red (high)
	(contour interval varies between panels). The graphs show solution for
	$l=0$. On a fine grid ($d \ll R_D$, high-resolution) the C-grid gives the best results.  
	
	\ESS
}

\section{Arakawa's 2D grids}
\subsection{}
\frame { \frametitle{Horizontally staggered grids} 

\BSS Need to choose finite difference
methods to solve the {\bf shallow-water equations} on the
plane tangent to the Earth's surface, {\em linearised} about a resting basic state:
\begin{eqnarray}
&&\frac{\partial \eta}{\partial t} + H \left(
\frac{\partial u}{\partial x}+\frac{\partial v}{\partial y}\right) = 0,
\\
&&\frac{\partial u}{\partial t} - f_0 v = - g\frac{\partial \eta}{\partial x}, \\
&&\frac{\partial v}{\partial t} + f_0 u = - g\frac{\partial \eta}{\partial y}, 
\end{eqnarray}
where $\eta$ is the surface elevation, $H$ is fluid depth, $(u,v)$ is
the depth-averaged horizontal velocity, $f$ is the Coriolis parameter
and $g$ is the gravitational acceleration. 

~~

Now discretise the problem using regular grids to describe $u$, $v$ and $\eta$ 

e.g., $\eta_{ij}=$ free surface elevation at the $i$th longitude and $j$th
latitude point.

~

Arakawa\footnote{\BTi Arakawa
A. (1966) Computational design for long-term numerical integration of
the equations of fluid motion: Two-dimensional incompressible
flow. Part I. \emph{Journal of Computational Physics}, \textbf{1},
119-143.\ETi} proposed a number of staggered finite difference
grids that can be used to solve the shallow water equations.  

\ESS }


\frame
{
\frametitle{Arakawa's A grid}
\BSS
Simplest and intuitive, but there is the danger of computational modes.
\begin{tabular}{lc}
\begin{minipage}[c]{0.6\textwidth}
\begin{eqnarray*}
\eta-grid \;\;\; \;\;\; \;\;\;
\frac{\partial u}{\partial x}\Big|_{ij} &\approx& \frac{u_{i+1,j}-u_{i-1,j}}{2 \Delta x} \\
\frac{\partial v}{\partial y}\Big|_{ij} &\approx& \frac{v_{i,j+1}-v_{i,j-1}}{2 \Delta y} \\
&& \\
u-grid \;\;\; \;\;\; \;\;\;
\frac{\partial \eta}{\partial x}\Big|_{ij} &\approx& \frac{\eta_{i+1,j}-\eta_{i-1,j}}{2 \Delta x} \\
f v\big|_{ij} &\approx& f v_{i,j} \\
&& \\
v-grid \;\;\; \;\;\; \;\;\;
\frac{\partial \eta}{\partial y}\Big|_{ij} &\approx& \frac{\eta_{i,j+1}-\eta_{i,j-1}}{2 \Delta y} \\
f u\big|_{ij} &\approx& f u_{i,j}
\end{eqnarray*}
\end{minipage}
&
\begin{minipage}[c]{0.4\textwidth}
\setlength{\unitlength}{1 cm}
\begin{picture}(4,4)
\arakawa
\put(0.93,0.93){$\bullet$} \put(1,1){\vector(1,0){0.5}} \put(1,1){\vector(0,1){0.5}}
\put(2.93,0.93){$\bullet$} \put(3,1){\vector(1,0){0.5}} \put(3,1){\vector(0,1){0.5}}
\put(0.93,2.93){$\bullet$} \put(1,3){\vector(1,0){0.5}} \put(1,3){\vector(0,1){0.5}}
\put(2.93,2.93){$\bullet$} \put(3,3){\vector(1,0){0.5}} \put(3,3){\vector(0,1){0.5}}
\put(1.1,0.7){$u_{ij}$,$v_{ij}$,}
\put(1.1,0.4){$\eta_{ij}$}
\put(3.1,3.6){$u_{i+1,j+1}$}
\put(3.1,3.4){$v_{i+1,j+1}$}
\put(3.1,3.2){$\eta_{i+1,j+1}$}
\put(3.1,0.7){$u_{i+1,j}$}
\put(3.1,0.5){$v_{i+1,j}$}
\put(3.1,0.3){$\eta_{i+1,j}$}
\end{picture}
\end{minipage}
\end{tabular}
\ESS
}


\frame
{
\frametitle{Arakawa's B grid}
\BSS
Stagger the variables, eliminate computational mode, but lots of interpolation...
\begin{tabular}{lc}
\begin{minipage}[c]{0.6\textwidth}
\begin{eqnarray*}
\frac{\partial u}{\partial x}\Big|_{ij} &\approx& \left(\frac{u_{i+1,j}+u_{i+1,j+1}}{2} - \frac{u_{i,j}+u_{i,j+1}}{2} \right)\frac{1}{\Delta x} \\
\frac{\partial v}{\partial y}\Big|_{ij} &\approx& \left(\frac{v_{i,j+1}+v_{i+1,j+1}}{2} - \frac{v_{i,j}+v_{i+1,j}}{2} \right)\frac{1}{\Delta y} \\
&& \\
\frac{\partial \eta}{\partial x}\Big|_{ij} &\approx& \left(\frac{\eta_{i,j-1}+\eta_{i,j}}{2} - \frac{\eta_{i-1,j-1}+\eta_{i-1,j}}{2} \right)\frac{1}{\Delta x} \\
f v\big|_{ij} &\approx& f v_{i,j} \\
&& \\
\frac{\partial \eta}{\partial y}\Big|_{ij} &\approx& \left(\frac{\eta_{i-1,j}+\eta_{i,j}}{2} - \frac{\eta_{i-1,j-1}+\eta_{i,j-1}}{2} \right)\frac{1}{\Delta y} \\
f u\big|_{ij} &\approx& f u_{i,j}
\end{eqnarray*}
\end{minipage}
&
\begin{minipage}[c]{0.4\textwidth}
\setlength{\unitlength}{1 cm}
\begin{picture}(4,4)
\arakawa
\put(1.93,1.93){$\bullet$}
\put(2.1,1.7){$\eta_{ij}$}
\put(1,1){\vector(1,0){0.5}} \put(1,1){\vector(0,1){0.5}}
\put(3,1){\vector(1,0){0.5}} \put(3,1){\vector(0,1){0.5}}
\put(1,3){\vector(1,0){0.5}} \put(1,3){\vector(0,1){0.5}}
\put(3,3){\vector(1,0){0.5}} \put(3,3){\vector(0,1){0.5}}
\put(1.1,0.7){$u_{ij}$,$v_{ij}$}
\put(3.1,3.4){$u_{i+1,j+1}$}
\put(3.1,3.2){$v_{i+1,j+1}$}
\put(3.1,0.7){$u_{i+1,j}$}
\put(3.1,0.5){$v_{i+1,j}$}
\end{picture}
\end{minipage}
\end{tabular}
\ESS
}


\frame
{
\frametitle{Arakawa's C grid}
\BSS
Stagger the variables, eliminate computational mode, interpolation only required for Coriolis terms!
\begin{tabular}{lc}
\begin{minipage}[c]{0.6\textwidth}
\begin{eqnarray*}
\eta-grid \;\;\; \;\;\; 
\frac{\partial u}{\partial x}\Big|_{ij} &\approx& \frac{u_{i+1,j}-u_{i,j}}{\Delta x} \\
\frac{\partial v}{\partial y}\Big|_{ij} &\approx& \frac{v_{i,j+1}-v_{i,j}}{\Delta y} \\
&& \\
u-grid \;\;\; \;\;\; 
\frac{\partial \eta}{\partial x}\Big|_{ij} &\approx& \frac{\eta_{i,j}-\eta_{i-1,j}}{\Delta x} \\
f v\big|_{ij} &\approx& f_j \frac{v_{i,j}+v_{i,j+1}+v_{i-1,j}+v_{i-1,j+1}}{4} \\
&& \\
v-grid \;\;\; \;\;\; 
\frac{\partial \eta}{\partial y}\Big|_{ij} &\approx& \frac{\eta_{i,j}-\eta_{i,j-1}}{\Delta y} \\
f u\big|_{ij} &\approx& f_j \frac{u_{i,j}+u_{i+1,j}+u_{i,j-1}+u_{i+1,j-1}}{4}
\end{eqnarray*}
\end{minipage}
&
\begin{minipage}[c]{0.4\textwidth}
\setlength{\unitlength}{1 cm}
\begin{picture}(4,4)
\arakawa
\put(1.93,1.93){$\bullet$}
\put(2.1,1.7){$\eta_{ij}$}
\put(2,0.75){\vector(0,1){0.5}}
\put(2,2.75){\vector(0,1){0.5}}
\put(0.75,2){\vector(1,0){0.5}}
\put(2.75,2){\vector(1,0){0.5}}
\put(1.1,1.7){$u_{ij}$}
\put(2.1,0.7){$v_{ij}$}
\put(3.1,1.7){$u_{i+1,~j}$}
\put(2.1,2.7){$v_{i,~j+1}$}
\end{picture}
\end{minipage}
\end{tabular}
\ESS

}

\frame
{
\frametitle{Arakawa's D grid}
\BSS
Idea similar to B grid, but we push out $u,v$ instead of $\eta$: stagger the variables, eliminate computational mode, but lots of interpolation...
\begin{tabular}{lc}
\begin{minipage}[c]{0.7\textwidth}
\begin{eqnarray*}
\frac{\partial u}{\partial x}\Big|_{ij} &\approx& \left(\frac{u_{i+1,j}+u_{i+1,j-1}}{2} - \frac{u_{i-1,j}+u_{i-1,j-1}}{2} \right)\frac{1}{2 \Delta x} \\
\frac{\partial v}{\partial y}\Big|_{ij} &\approx& \left(\frac{v_{i,j+1}+v_{i-1,j+1}}{2} - \frac{v_{i,j-1}+v_{i-1,j-1}}{2} \right)\frac{1}{2 \Delta y} \\
&& \\
\frac{\partial \eta}{\partial x}\Big|_{ij} &\approx& \left(\frac{\eta_{i+1,j}+\eta_{i+1,j+1}}{2} - \frac{\eta_{i-1,j}+\eta_{i-1,j+1}}{2} \right)\frac{1}{2 \Delta x} \\
f v\big|_{ij} &\approx& f_j \frac{v_{i,j}+v_{i,j+1}+v_{i-1,j}+v_{i-1,j+1}}{4} \\
&& \\
\frac{\partial \eta}{\partial y}\Big|_{ij} &\approx& \left(\frac{\eta_{i,j+1}+\eta_{i+1,j+1}}{2} - \frac{\eta_{i,j-1}+\eta_{i+1,j-1}}{2} \right)\frac{1}{2 \Delta y} \\
f u\big|_{ij} &\approx& f_j \frac{u_{i,j}+u_{i+1,j}+u_{i,j-1}+u_{i+1,j-1}}{4}
\end{eqnarray*}
\end{minipage}
&
\begin{minipage}[c]{0.3\textwidth}
\setlength{\unitlength}{1 cm}
\begin{picture}(4,4)
\arakawa
\put(0.93,0.93){$\bullet$}
\put(2.93,0.93){$\bullet$}
\put(0.93,2.93){$\bullet$}
\put(2.93,2.93){$\bullet$}
\put(2,0.75){\vector(0,1){0.5}}
\put(2,2.75){\vector(0,1){0.5}}
\put(0.75,2){\vector(1,0){0.5}}
\put(2.75,2){\vector(1,0){0.5}}
\put(1.1,1.7){$u_{ij}$}
\put(2.1,0.7){$v_{ij}$}
\put(1.1,0.7){$\eta_{ij}$}
\end{picture}
\end{minipage}
\end{tabular}
\ESS

}

\frame
{
	\frametitle{Arakawa's Sm\"org{\aa}sbord}
	\BSS
\begin{figure}[h]
	\centering
	\includegraphics[width=0.8\textwidth]{/Users/vidale/Documents/Presentations/"SWE 2D grid stagger types".pdf}
	\caption{Many of Arakawa's grid staggers}
	\label{fig:Arakawa_all}
\end{figure}

\ESS

}

\section{Excursus: SWEs in 1D}
\subsection{}
\frame { \frametitle{Excursus: 1D SWEs help to understand grid stagger} 
	\BSS 
	\includegraphics[trim={0cm 0.5cm 0cm 2.cm},clip,width=11.cm]{/Users/vidale/Projects/Teaching/MTMW14/"2016 module"/Adv_2_SWEs_1.pdf}
	\label{fig:Spherical-PEs1}
	\setbeamerfont{caption name}{size=\tiny}
	%\caption{\label{fig:blue_rectangle} \tiny PEs on the Sphere}
	\ESS }

\frame { \frametitle{Excursus: resolving waves on a grid} 
	\BSS 
	\includegraphics[trim={0cm 0.cm 0cm 0.cm},clip,width=9.cm]{/Users/vidale/Projects/Teaching/MTMW14/"2019 module/Waves resolved on a Grid".png}
	\label{fig:Spherical-PEs2}
	\setbeamerfont{caption name}{size=\tiny}
	%\caption{\label{fig:blue_rectangle} \tiny PEs on the Sphere}
	\ESS }

\frame { \frametitle{Excursus: 1D SWEs help to understand grid stagger} 
	\BSS 
	\includegraphics[trim={0cm 0.5cm 0cm 2.cm},clip,width=11.cm]{/Users/vidale/Projects/Teaching/MTMW14/"2016 module"/Adv_2_SWEs_2.pdf}
	\label{fig:Spherical-PEs2}
	\setbeamerfont{caption name}{size=\tiny}
	%\caption{\label{fig:blue_rectangle} \tiny PEs on the Sphere}
	\ESS }

\frame { \frametitle{Excursus: 1D SWEs help to understand grid stagger} 
	\BSS 
	\includegraphics[trim={0cm 0.5cm 0cm 2.cm},clip,width=11.cm]{/Users/vidale/Projects/Teaching/MTMW14/"2016 module"/Adv_2_SWEs_3.pdf}
	\label{fig:Spherical-PEs2}
	\setbeamerfont{caption name}{size=\tiny}
	%\caption{\label{fig:blue_rectangle} \tiny PEs on the Sphere}
	\ESS }

\frame { \frametitle{Excursus: 1D SWEs help to understand grid stagger} 
	\BSS 
	\includegraphics[trim={0cm 0.5cm 0cm 1.cm},clip,width=11.cm]{/Users/vidale/Projects/Teaching/MTMW14/"2016 module"/Adv_2_SWEs_4.pdf}
	\label{fig:Spherical-PEs2}
	\setbeamerfont{caption name}{size=\tiny}
	%\caption{\label{fig:blue_rectangle} \tiny PEs on the Sphere}
	\ESS }

\frame { \frametitle{Exercise: analytical versus numerical solutions} 
	\BSS 
	Start from the two shallow water equations on the previous page. Impose the wave solution proposed on the slide, for $u_j$ and $h_j$: 
	\begin{equation}
		(u_j,h_j) \sim e^{i(k x_j -\omega t)} = e^{i(k j \Delta x -\omega t)}
	\end{equation}. 
	\begin{enumerate}
		\item Substitute into the centred numerical expressions (in time and space) and show that you end up with an expression that contains a $\sin$ function.
		\item How does the numerical solution compare to the analytical solution in terms of a) amplitude and b) phase?
		\item Plot a wavenumber versus frequency plot for the analytical and numerical solution
		\item Expand your thinking to Numerical Weather Prediction models that might be using such schemes: discuss what operational problems we may encounter in the presence of a $\sin$ in such numerical predictions.
	\end{enumerate}
	\ESS }

\end{document}